% 书籍配图LaTeX引用汇总
% 使用方法：将所需图片代码复制到LaTeX文档中

\usepackage{graphicx}
\usepackage{float}
\usepackage{caption}
\usepackage{subcaption}

% ============================================================
% 第3章 - DSF/BAG架构配图
% ============================================================

% 图3.1 DSF双域架构示意图
\begin{figure}[H]
    \centering
    \includegraphics[width=0.95\textwidth]{chapter3/dsf-dual-domain.png}
    \caption{DSF双域架构示意图。以太网域提供传统网络服务，交换结构域提供超低延迟的调度fabric，两者协同工作。}
    \label{fig:dsf-dual-domain}
\end{figure}

% 图3.2 DSF核心组件架构
\begin{figure}[H]
    \centering
    \includegraphics[width=0.95\textwidth]{chapter3/dsf-components.png}
    \caption{DSF核心组件架构。IN（RDSW）负责服务器接入，FN（FDSW）负责高速交换，通过VOQ调度实现无阻塞传输。}
    \label{fig:dsf-components}
\end{figure}

% 图3.3 BAG跨区域互联架构
\begin{figure}[H]
    \centering
    \includegraphics[width=0.95\textwidth]{chapter3/bag-topology.png}
    \caption{BAG跨区域互联架构。多个数据中心的DSF通过BAG超级脊层互联，形成大规模AI训练网络。}
    \label{fig:bag-topology}
\end{figure}

% 图3.4 包喷洒技术示意图
\begin{figure}[H]
    \centering
    \includegraphics[width=0.95\textwidth]{chapter3/packet-spraying.png}
    \caption{包喷洒技术示意图。通过将数据流分散到多条路径传输，提高网络利用率和避免拥塞，接收端通过VOQ缓冲区重排序。}
    \label{fig:packet-spraying}
\end{figure}

% ============================================================
% 第4章 - P4/DPU/CXL配图
% ============================================================

% 图4.1 P4可编程数据平面架构
\begin{figure}[H]
    \centering
    \includegraphics[width=0.95\textwidth]{chapter4/p4-programmable-plane.png}
    \caption{P4可编程数据平面架构。通过P4语言定义数据包处理逻辑，编译后部署到可编程交换机或智能网卡。}
    \label{fig:p4-architecture}
\end{figure}

% 图4.2 DPU架构与异构计算协作
\begin{figure}[H]
    \centering
    \includegraphics[width=0.95\textwidth]{chapter4/dpu-architecture.png}
    \caption{DPU架构与异构计算协作。DPU承担网络、存储、安全等基础设施任务，释放CPU和GPU专注于业务计算。}
    \label{fig:dpu-architecture}
\end{figure}

% 图4.3 CXL内存扩展架构
\begin{figure}[H]
    \centering
    \includegraphics[width=0.95\textwidth]{chapter4/cxl-memory-expansion.png}
    \caption{CXL内存扩展架构。通过CXL协议实现内存池化、设备缓存一致性和异构内存扩展，突破单节点内存限制。}
    \label{fig:cxl-architecture}
\end{figure}

% 图4.4 eBPF+Cilium数据路径
\begin{figure}[H]
    \centering
    \includegraphics[width=0.95\textwidth]{chapter4/ebpf-cilium-datapath.png}
    \caption{eBPF+Cilium数据路径架构。利用eBPF在内核关键路径插入自定义逻辑，实现高性能的网络策略执行和可观测性。}
    \label{fig:ebpf-cilium}
\end{figure}

% ============================================================
% 第5章 - RoCEv2/RDMA配图
% ============================================================

% 图5.1 RoCEv2与InfiniBand对比
\begin{figure}[H]
    \centering
    \includegraphics[width=0.95\textwidth]{chapter5/rocev2-vs-infiniband.png}
    \caption{RoCEv2与InfiniBand协议栈对比。两者共享Verbs API，但在网络层实现上有差异，RoCEv2更适合大规模数据中心部署。}
    \label{fig:roce-vs-ib}
\end{figure}

% 图5.2 DCQCN拥塞控制流程
\begin{figure}[H]
    \centering
    \includegraphics[width=0.95\textwidth]{chapter5/dcqcn-congestion-control.png}
    \caption{DCQCN拥塞控制流程。交换机通过ECN标记拥塞，接收端生成CNP通知发送端，发送端采用量化的速率调整算法降低发送速率。}
    \label{fig:dcqcn-flow}
\end{figure}

% 图5.3 RDMA数据传输路径
\begin{figure}[H]
    \centering
    \includegraphics[width=0.95\textwidth]{chapter5/rdma-data-path.png}
    \caption{RDMA数据传输路径对比。传统TCP/IP需要多次内存拷贝和内核介入，RDMA通过内核旁路实现零拷贝传输，大幅降低延迟和CPU开销。}
    \label{fig:rdma-path}
\end{figure}

% 图5.4 负载均衡技术演进
\begin{figure}[H]
    \centering
    \includegraphics[width=0.95\textwidth]{chapter5/load-balancing-evolution.png}
    \caption{数据中心负载均衡技术演进时间线。从基于流的ECMP到包级喷洒，负载均衡粒度不断细化，适应AI训练等大带宽需求场景。}
    \label{fig:lb-evolution}
\end{figure}

% ============================================================
% 第6章 - AI训练网络配图
% ============================================================

% 图6.1 BAG跨区域互联拓扑
\begin{figure}[H]
    \centering
    \includegraphics[width=0.95\textwidth]{chapter6/bag-interconnect.png}
    \caption{BAG跨区域互联拓扑。通过BAG超级脊层连接多个区域的数据中心，实现跨地域的大规模AI训练集群。}
    \label{fig:bag-interconnect}
\end{figure}

% 图6.2 DSF三级扩展架构
\begin{figure}[H]
    \centering
    \includegraphics[width=0.95\textwidth]{chapter6/dsf-three-tier.png}
    \caption{DSF三级扩展架构。L1实现机架内高速互联，L2实现集群内全连接，L3通过BAG实现跨集群扩展，支持百万级GPU规模。}
    \label{fig:dsf-tiers}
\end{figure}

% 图6.3 AI训练集群网络架构全景
\begin{figure}[H]
    \centering
    \includegraphics[width=0.95\textwidth]{chapter6/ai-cluster-overview.png}
    \caption{AI训练集群网络架构全景。前端网络负责管理和存储，后端网络负责训练计算，通过分离设计实现最佳性能。}
    \label{fig:ai-cluster-overview}
\end{figure}

% 图6.4 性能优化工具链
\begin{figure}[H]
    \centering
    \includegraphics[width=0.95\textwidth]{chapter6/performance-toolchain.png}
    \caption{AI训练网络性能优化工具链。从指标采集到智能分析，再到自动化优化，形成完整的闭环性能管理系统。}
    \label{fig:perf-toolchain}
\end{figure}

% ============================================================
% 第7章 - 绿色计算配图
% ============================================================

% 图7.1 数据中心PUE优化示意图
\begin{figure}[H]
    \centering
    \includegraphics[width=0.9\textwidth]{chapter7/pue-optimization.png}
    \caption{数据中心PUE优化示意图。PUE越接近1.0表示能效越高，主要通过优化冷却系统、采用液冷技术、利用自然冷却和可再生能源来实现。}
    \label{fig:pue-optimization}
\end{figure}

% 图7.2 液冷技术架构对比图
\begin{figure}[H]
    \centering
    \includegraphics[width=0.95\textwidth]{chapter7/liquid-cooling-architecture.png}
    \caption{液冷技术架构对比图。冷板式液冷适合渐进式部署，浸没式液冷提供最高散热效率但需要专门设计。}
    \label{fig:liquid-cooling}
\end{figure}

% 图7.3 可再生能源应用架构
\begin{figure}[H]
    \centering
    \includegraphics[width=0.95\textwidth]{chapter7/renewable-energy.png}
    \caption{数据中心可再生能源应用架构。通过多源互补发电、储能缓冲和电网互联，实现高比例可再生能源供电。}
    \label{fig:renewable-energy}
\end{figure}

% 图7.4 Meta可持续发展时间线
\begin{figure}[H]
    \centering
    \includegraphics[width=0.95\textwidth]{chapter7/meta-sustainability-timeline.png}
    \caption{Meta数据中心可持续发展时间线（2011-2024）。从Open Compute Project到100\%可再生能源，再到AI时代的可持续基础设施。}
    \label{fig:meta-timeline}
\end{figure}

% ============================================================
% 图片引用示例
% ============================================================

% 在正文中引用图片：
% 如图~\ref{fig:dsf-dual-domain}所示，DSF采用双域架构设计...
% 根据图~\ref{fig:rdma-path}的对比，RDMA相比传统TCP/IP具有显著优势...

% 多个子图示例：
% \begin{figure}[H]
%     \centering
%     \begin{subfigure}[b]{0.48\textwidth}
%         \centering
%         \includegraphics[width=\textwidth]{chapter5/rocev2-stack.png}
%         \caption{RoCEv2协议栈}
%         \label{fig:roce-stack}
%     \end{subfigure}
%     \hfill
%     \begin{subfigure}[b]{0.48\textwidth}
%         \centering
%         \includegraphics[width=\textwidth]{chapter5/infiniband-stack.png}
%         \caption{InfiniBand协议栈}
%         \label{fig:ib-stack}
%     \end{subfigure}
%     \caption{RoCEv2与InfiniBand协议栈对比}
%     \label{fig:protocol-comparison}
% \end{figure}
